% Content for the 90-minute seminar "AI in the Flipped Classroom" (hands-on; Sharad Pawar Inspire
% Fellowship in Education, 31 Oct 2026). Used by Main_Seminar_AI_For_Educators_Flipped_Presentation.tex
% and _CheatSheet.tex. The full 1-hour seminar is Main_Seminar_AI_For_Educators_* (no flipped module).
% The AI introduction and CRAFT come from the _short siblings of the full seminar's files
% (ai_educators_content_short, llm_promptengg_educators_short): same frames, a few commented out.
% Run of show (90 min): AI and education 18, CRAFT 10, Flipped 5, five Exercises 45 (9 min each),
% Checks 7, Takeaways, promo and questions 5.
% Every listing is the last element of its frame. Prompts use language={} (plain text, no colouring).
% CheatSheet-safe: no floats; the one table is wrapped in \resizebox.

% Sections HowAI, State, TaskvJob (condensed): see ai_educators_content_short.tex
% SHORT sibling of ai_educators_content.tex (created 2026-10-24).
% Full copy with the frames not needed in the 90-minute flipped seminar commented out, never deleted.
% Mirror any frame edit between the two files; restore a frame by removing the % prefixes.
% Merged 2026: folds ai_educators_howaiworks/_stateofplay/_taskvsjob/_programstructure
% into one file (they were 4 small standalone files, 15-72 lines each). Section boundaries
% are marked inline with \section so this single \input still produces 4 separate sections
% in the deck, exactly as before the merge.

\section[WhatAI]{What Is AI?}

%%%%%%%%%%%%%%%%%%%%%%%%%%%%%%%%%%%%%%%%%%%%%%%%%%%%%%%%%%%
% [upgrade-deck: moved 2026-10-24 | reorder | AI/ML/DL big picture now opens the deck; same move in ai_educators_content.tex]
\begin{frame}[fragile]\frametitle{The Big Picture: AI, ML, and Deep Learning}
\begin{center}
\includegraphics[width=0.55\linewidth,keepaspectratio]{ai1}
\end{center}
\begin{itemize}
\item \textbf{AI}: any machine doing something that looks intelligent
\item \textbf{Machine Learning (ML)}: a machine that learns the pattern from data, instead of being told the rule directly
\item \textbf{Deep Learning (DL)}: ML using many-layered neural networks, powerful enough to learn something as complex as language
\end{itemize}
{\tiny (Ref: https://blogs.nvidia.com/blog/2016/07/29/whats-difference-artificial-intelligence-machine-learning-deep-learning-ai/)}
\end{frame}

\section[HowAI]{How Generative AI Actually Works}

%%%%%%%%%%%%%%%%%%%%%%%%%%%%%%%%%%%%%%%%%%%%%%%%%%%%%%%%%%%%%%%%%%%%%%%%%%%%%%%%%%
\begin{frame}[fragile]\frametitle{}
\begin{center}
{\Large How Generative AI Actually Works}

\medskip
{\small (Just enough to demystify it, not a computer science course)}
\end{center}
\end{frame}

%%%%%%%%%%%%%%%%%%%%%%%%%%%%%%%%%%%%%%%%%%%%%%%%%%%%%%%%%%%
\begin{frame}[fragile]\frametitle{What Is a Language Model?}
\begin{itemize}
\item While typing an SMS, have you noticed it suggests the next word?
\item Those suggestions are not random: the phone has learned what word usually comes next from millions of past sentences.
\item Do this once, word by word, and you get a suggestion. Do it again and again, and you get a full sentence.
\item That is a \textbf{Language Model}, predicting the next word. Chained together, it writes whole paragraphs: a \textbf{Generative Model}.
\end{itemize}
\begin{center}
\includegraphics[width=0.4\linewidth,keepaspectratio]{chatgpt34}
\end{center}
\end{frame}

%%%%%%%%%%%%%%%%%%%%%%%%%%%%%%%%%%%%%%%%%%%%%%%%%%%%%%%%%%%
\begin{frame}[fragile]\frametitle{Why Are They Called ``Large''?}
\begin{itemize}
\item ChatGPT works on exactly this idea, just at an enormous scale: trained on a large fraction of everything ever written on the internet.
\item That is why the phrase is \textbf{Large Language Model (LLM)}: large training data, large number of internal parameters.
\item More data and more parameters mean better predictions, but the underlying idea never changes: predict the next word, one at a time.
\end{itemize}
\begin{center}
\includegraphics[width=0.5\linewidth,keepaspectratio]{chatgpt31}
\end{center}
\end{frame}

\section[State]{Where AI in Education Stands}

%%%%%%%%%%%%%%%%%%%%%%%%%%%%%%%%%%%%%%%%%%%%%%%%%%%%%%%%%%%%%%%%%%%%%%%%%%%%%%%%%%
\begin{frame}[fragile]\frametitle{}
\begin{center}
{\Large Where AI in Education Actually Stands}
\end{center}
\end{frame}

%%%%%%%%%%%%%%%%%%%%%%%%%%%%%%%%%%%%%%%%%%%%%%%%%%%
\begin{frame}[fragile]\frametitle{What Teachers Are Actually Using It For}
\begin{itemize}
\item Teacher AI adoption has grown fast worldwide over the last two years, surveys everywhere point the same direction.
\item The most common uses are unglamorous but time-consuming: gathering research and content, drafting lesson plans, summarizing information, generating classroom materials.
\item \textbf{Notice what all four have in common:} none of them is ``teach the class'' or ``know the student.''
\end{itemize}
\end{frame}

%%%%%%%%%%%%%%%%%%%%%%%%%%%%%%%%%%%%%%%%%%%%%%%%%%%
\begin{frame}[fragile]\frametitle{Opportunities AI Opens Up}
\begin{itemize}
\item \textbf{Differentiated learning, without extra prep:} the same worksheet re-explained at three different difficulty levels, generated in the time it used to take to write one.
\item \textbf{Patient, always-on practice for struggling students:} a student who needs the same concept explained five different ways gets it without feeling like a burden on the teacher or the class.
\item \textbf{A bridge for the language gap:} explaining or translating a difficult English textbook concept into Marathi (or any regional language) on the spot, so understanding does not wait for a bilingual teacher to be free.
\item \textbf{Reclaimed time:} the hours saved on drafting and admin (previous slide) are hours that can go back into actually noticing and helping individual students.
\end{itemize}
\end{frame}

\section[TaskvJob]{Will AI Replace Teachers?}

%%%%%%%%%%%%%%%%%%%%%%%%%%%%%%%%%%%%%%%%%%%%%%%%%%%%%%%%%%%%%%%%%%%%%%%%%%%%%%%%%%
% [upgrade-deck: commented out 2026-10-24 | short sibling | left out of the 90-minute flipped seminar (full version keeps it)]
% \begin{frame}[fragile]\frametitle{}
% \begin{center}
% {\Large Will AI Replace Teachers?}
% \end{center}
% \end{frame}

%%%%%%%%%%%%%%%%%%%%%%%%%%%%%%%%%%%%%%%%%%%%%%%%%%%
% [upgrade-deck: restored 2026-10-24 | short sibling | Hinton box brought back at the author's request]
\begin{frame}[fragile]\frametitle{2016: A Very Confident Prediction}
\begin{itemize}
\item Every major AI advance since the 1980s has come with the same headline: ``this job is over.'' Radiology got there first, and loudest, a full decade before ChatGPT existed, so what actually happened to it is the single most useful data point we have for thinking about teaching.
\item Geoffrey Hinton, the ``Godfather of AI'' and a Nobel laureate, told a room of doctors in 2016:
\end{itemize}
\begin{alertblock}{Hinton, 2016}
``We should stop training radiologists now. It's just completely obvious that within five years, deep learning is going to do better than radiologists.''
\end{alertblock}
\begin{itemize}
\item He later said, to the New York Times, that he had spoken too broadly.
\end{itemize}
{\tiny (Ref: New York Times, revisited in \emph{radiologybusiness.com}, 2026.)}
\end{frame}

%%%%%%%%%%%%%%%%%%%%%%%%%%%%%%%%%%%%%%%%%%%%%%%%%%%
% [upgrade-deck: restored 2026-10-24 | short sibling | the Hinton story needs its outcome slide]
\begin{frame}[fragile]\frametitle{What Actually Happened}
\begin{itemize}
\item Radiologists did not disappear. Demand for them grew, and it is still growing.
\item AI got very good at one task inside the job: reading a scan.
\item Everything else, consulting with physicians, explaining findings to a worried patient, deciding what to biopsy, stayed entirely human.
\item \textbf{The lesson:} a job survives as long as most of its tasks resist automation, even when one high-profile task falls.
\end{itemize}
\end{frame}

%%%%%%%%%%%%%%%%%%%%%%%%%%%%%%%%%%%%%%%%%%%%%%%%%%%
\begin{frame}[fragile]\frametitle{The Same Logic Applies to Teaching}
\begin{itemize}
\item ``Teacher'' is also a bundle of tasks, not one thing.
\item \textbf{Highly automatable:} drafting a first-pass lesson outline, generating quiz questions, summarizing 50 feedback forms, writing a first draft of report comments.
\item \textbf{Not automatable any time soon:} noticing a student is struggling before they say so, managing a room of 30 restless 12-year-olds, a difficult conversation with a parent, being the adult a kid trusts.
\item AI is quietly eating the first list. Nobody is close to touching the second.
\end{itemize}
\end{frame}

%%%%%%%%%%%%%%%%%%%%%%%%%%%%%%%%%%%%%%%%%%%%%%%%%%%
\begin{frame}[fragile]\frametitle{Quick Check: Task or Job?}
\begin{itemize}
\item Nvidia's CEO Jensen Huang put it this way in 2026: \emph{``AI kills tasks, not jobs.''}
\item For each of these, is AI automating a \textbf{task} within teaching, or threatening the \textbf{job} of teaching itself?
\begin{itemize}
\item Grading 40 multiple-choice quizzes
\item Resolving a conflict between two students at recess
\item Drafting a first pass at a rubric
\item Deciding whether a struggling student needs extra support
\end{itemize}
\end{itemize}
\end{frame}

%%%%%%%%%%%%%%%%%%%%%%%%%%%%%%%%%%%%%%%%%%%%%%%%%%%
\begin{frame}[fragile]\frametitle{Quick Check: Answer}
\begin{itemize}
\item Grading MCQs, drafting a rubric: \textbf{tasks}. Mechanical, pattern-based, already being automated, and that is fine.
\item Resolving a peer conflict, judging who needs support: \textbf{the job}. These need judgment, context, and a relationship built over months, exactly what nobody has figured out how to automate.
\item \textbf{The reframe:} the question is never ``will AI replace teachers,'' it is ``which specific tasks on my plate can I hand off this term.''
\end{itemize}
\end{frame}

% [upgrade-deck: commented out 2026-10-24 | short sibling | left out of the 90-minute flipped seminar (full version keeps it)]
% \section[Programs]{How These Programs Are Structured}

%%%%%%%%%%%%%%%%%%%%%%%%%%%%%%%%%%%%%%%%%%%%%%%%%%%%%%%%%%%%%%%%%%%%%%%%%%%%%%%%%%
% [upgrade-deck: commented out 2026-10-24 | short sibling | left out of the 90-minute flipped seminar (full version keeps it)]
% \begin{frame}[fragile]\frametitle{}
% \begin{center}
% {\Large How ``AI for Teachers'' Programs Are Actually Structured}
% \end{center}
% \end{frame}

%%%%%%%%%%%%%%%%%%%%%%%%%%%%%%%%%%%%%%%%%%%%%%%%%%%
% [upgrade-deck: commented out 2026-10-24 | short sibling | left out of the 90-minute flipped seminar (full version keeps it)]
% \begin{frame}[fragile]\frametitle{The Four Pieces, Every Time}
% \begin{itemize}
% \item Whatever the branding, every serious education-AI program is built from the same four pieces.
% \item \textbf{A skills library:} ready-made prompts and workflows for recurring teacher tasks, lesson planning, rubric drafting, differentiated explanations, so nobody starts from a blank prompt box.
% \item \textbf{Standards grounding:} lesson plans are checked against your actual curriculum, not the model's generic training knowledge.
% \item \textbf{Task automation:} tools that run a recurring job on their own, e.g. a daily review of submitted exit tickets, without you re-prompting every time.
% \item \textbf{Privacy guardrails:} clear, explicit rules about what happens to student data.
% \end{itemize}
% \end{frame}

%%%%%%%%%%%%%%%%%%%%%%%%%%%%%%%%%%%%%%%%%%%%%%%%%%%
% [upgrade-deck: commented out 2026-10-24 | short sibling | left out of the 90-minute flipped seminar (full version keeps it)]
% \begin{frame}[fragile]\frametitle{What to Look For in Any Tool Your School Considers}
% \begin{itemize}
% \item Does it know your curriculum, or is it just guessing from general training knowledge?
% \item Can it run on autopilot for the boring, repetitive parts, or does every use start from scratch?
% \item Does it say clearly, in plain language, what happens to student data?
% \item If the answer to any of these is unclear, ask before adopting it school-wide.
% \end{itemize}
% \end{frame}


\section[Limits]{What These Tools Do Well and Badly}

%%%%%%%%%%%%%%%%%%%%%%%%%%%%%%%%%%%%%%%%%%%%%%%%%%%
\begin{frame}[fragile]\frametitle{What These Tools Do Well and Badly}
\begin{itemize}
\item \textbf{Strong:} first drafts, rewriting at another level, structure, examples, questions, brainstorming.
\item \textbf{Weak:} exact facts, numbers, references, recent events, and your own class.
\item \textbf{Hallucination:} the tool can state a wrong answer in a confident tone.
\item So you use it for the draft. You check and decide.
\end{itemize}
\end{frame}

\section[Prompts]{Prompts with CRAFT}
% Divider plus the six CRAFT frames: see llm_promptengg_educators_short.tex
% SHORT sibling of llm_promptengg_educators.tex (created 2026-10-24).
% Full copy with the frames not needed in the 90-minute flipped seminar commented out, never deleted.
% Mirror any frame edit between the two files; restore a frame by removing the % prefixes.
%%%%%%%%%%%%%%%%%%%%%%%%%%%%%%%%%%%%%%%%%%%%%%%%%%%%%%%%%%%%%%%%%%%%%%%%%%%%%%%%%%
\begin{frame}[fragile]\frametitle{}
\begin{center}
{\Large Prompt Engineering for Educators}

\end{center}
\end{frame}

%%%%%%%%%%%%%%%%%%%%%%%%%%%%%%%%%%%%%%%%%%%%%%%%%%%%%%%%%%%
% [upgrade-deck: added 2026-10-24 | intro to prompts | same two frames in llm_promptengg_educators.tex]
\begin{frame}[fragile]\frametitle{What Is a Prompt?}
\begin{itemize}
\item A \textbf{prompt} is the text you type to the AI: a question, an instruction, or both.
\item The model reads your prompt and continues it, one word at a time, just like the SMS suggestion.
\item So your words steer the answer: more specific text gives a more specific result.
\item A prompt can hold a role for the AI, background, an example, and the text to work on.
\item \textbf{Prompt engineering} is writing prompts that give usable answers.
\end{itemize}
\end{frame}

%%%%%%%%%%%%%%%%%%%%%%%%%%%%%%%%%%%%%%%%%%%%%%%%%%%
\begin{frame}[fragile]\frametitle{Same Task, Different Prompts}
\begin{itemize}
\item Same AI, same topic: a vague prompt gives a generic answer, a specific one gives a usable answer.
\item A chat remembers what came before. Reply to refine the answer instead of starting again.
\end{itemize}
\begin{lstlisting}[language={}]
Vague:      Explain photosynthesis.

Specific:   Explain photosynthesis to a Grade 5 student
            in 5 sentences, with one everyday example.

Follow-up:  Make it shorter, and add one question to
            check understanding.
\end{lstlisting}
\end{frame}


%%%%%%%%%%%%%%%%%%%%%%%%%%%%%%%%%%%%%%%%%%%%%%%%%%%%%%%%%%%
% [upgrade-deck: commented out 2026-10-24 | short sibling | left out of the 90-minute flipped seminar (full version keeps it)]
% \begin{frame}[fragile]\frametitle{Why This Matters for You}
%
% \begin{columns}
%     \begin{column}[T]{0.55\linewidth}
%     You do not need to be a programmer to benefit from Prompt Engineering.
%
%     \medskip
%     \textbf{These are real tasks you do every week:}
%     \begin{itemize}
%         \item Planning lessons or training sessions
%         \item Giving feedback on student or colleague work
%         \item Organising your own thoughts before a meeting
%     \end{itemize}
%     \end{column}
%     \begin{column}[T]{0.45\linewidth}
%     \begin{center}
%     \renewcommand{\arraystretch}{1.4}
%     \begin{tabular}{|l|c|}
%     \hline
%     \textbf{Task} & \textbf{Time Saved} \\
%     \hline
%     Lesson plan & 60\% \\
%     Feedback writing & 75\% \\
%     \hline
%     \end{tabular}
%     \end{center}
%     \medskip
%     {\tiny Estimates based on typical knowledge worker time studies.}
%
% 	    \textbf{A well-crafted prompt can cut each of these tasks from 30 minutes to 3 minutes.}
%     \end{column}
% \end{columns}
%
% \end{frame}


%%%%%%%%%%%%%%%%%%%%%%%%%%%%%%%%%%%%%%%%%%%%%%%%%%%%%%%%%%%
% [upgrade-deck: commented out 2026-10-24 | short sibling | left out of the 90-minute flipped seminar (full version keeps it)]
% \begin{frame}[fragile]\frametitle{Creating a Lesson Plan (For Teachers)}
%
%     \begin{itemize}
%         \item \textbf{Situation:} You need to plan a 45-minute lesson on a topic for a specific age group, today.
%         \item \textbf{Result:} A complete, structured, ready-to-use lesson plan in under 30 seconds. Adapt the output to your classroom's specific needs.
%     \end{itemize}
%
% \begin{lstlisting}[language=HTML]
% Prompt:
% Create a detailed 45-minute lesson plan for:
% - Topic: Introduction to Fractions
% - Grade: 5th standard (age 10-11)
% - Prior knowledge: students know division basics
% - Learning objective: understand what a fraction means 
%   and compare simple fractions (1/2, 1/4, 3/4)
%
% Include:
% - 5-minute warm-up activity (no materials needed)
% - Main concept explanation with a real-life analogy
% - A hands-on activity using items in a classroom
% - 5 practice problems (easy to hard)
% - A 2-minute closing question to check understanding
% - One common misconception to watch out for
% \end{lstlisting}
%
%
%
% \end{frame}


%%%%%%%%%%%%%%%%%%%%%%%%%%%%%%%%%%%%%%%%%%%%%%%%%%%%%%%%%%%
% [upgrade-deck: commented out 2026-10-24 | short sibling | left out of the 90-minute flipped seminar (full version keeps it)]
% \begin{frame}[fragile]\frametitle{Generating Quiz and Exam Questions}
%
%     \begin{itemize}
%         \item \textbf{Situation:} You need to create 20 questions for an assessment. Writing good questions is time-consuming, especially writing good \emph{wrong} answers for multiple-choice.
%         \item \textbf{Why ``plausible wrong options'' matters:} Vague wrong answers make questions trivially easy. Asking the AI to use common misconceptions creates genuinely diagnostic questions.
%     \end{itemize}
%
% \begin{lstlisting}[language=HTML]
% Prompt:
% Create 10 multiple-choice questions on the topic: 
% "The Water Cycle" for Grade 7 students.
%
% For each question:
% - Write one clear question
% - Provide 4 options (A, B, C, D)
% - Make sure the 3 wrong options are plausible 
%   (common misconceptions, not obviously silly)
% - Mark the correct answer
% - Add a one-line explanation of why it is correct
%
% Difficulty: mix of easy (4), medium (4), hard (2).
% \end{lstlisting}
%
%
%
% \end{frame}


%%%%%%%%%%%%%%%%%%%%%%%%%%%%%%%%%%%%%%%%%%%%%%%%%%%%%%%%%%%
% [upgrade-deck: commented out 2026-10-24 | short sibling | left out of the 90-minute flipped seminar (full version keeps it)]
% \begin{frame}[fragile]\frametitle{Explaining Difficult Concepts at Different Levels}
%
%     \begin{itemize}
%         \item \textbf{Situation:} You need to explain the same concept to different audiences, a student, a parent, and the school principal.
%         \item \textbf{Why this is powerful:} You write the concept once; the AI adapts the vocabulary, analogies, and tone for each audience automatically.
%         \item \textbf{Office use:} Same technique for explaining a technical issue to your manager vs. to an IT colleague vs. to a client.
%     \end{itemize}
%
% \begin{lstlisting}[language=HTML]
% Prompt:
% Explain the concept of "compound interest" three times:
% 1. To a 12-year-old student using a pocket money analogy
% 2. To a parent at a school financial literacy evening, 
%    practical, relatable, no jargon
% 3. To a school board member in one precise paragraph 
%    for an annual report
%
% Keep each explanation self-contained and appropriate 
% for the audience.
% \end{lstlisting}
%
%
% \end{frame}

% Trimmed from the live deck (Aug 2026) to fit the merged 1-hour "AI for
% Educators" talk's frame budget; kept on disk, not deleted, for possible
% future re-inclusion.
% %%%%%%%%%%%%%%%%%%%%%%%%%%%%%%%%%%%%%%%%%%%%%%%%%%%%%%%%%%%
% \begin{frame}[fragile]\frametitle{Creating Training and Onboarding Materials}
%
%     \begin{itemize}
%         \item \textbf{Situation:} A new employee or student is joining. You need to create a guide, checklist, or FAQ quickly.
%         \item \textbf{Time saved:} What normally takes a senior staff member an afternoon of writing now takes 2 minutes to draft, 10 minutes to review and customise.
%     \end{itemize}
%
% \begin{lstlisting}[language=HTML]
% Prompt:
% Create a "First Week Onboarding Checklist" for a new
% primary school teacher joining mid-year.
%
% Include:
% - Day 1: Essential admin tasks (HR forms, ID card, etc.)
% - Day 2-3: Classroom setup and resource access
% - Day 4-5: Meeting key people (checklist format)
% - One-page FAQ answering the 8 most common questions
%   a new teacher asks (e.g., how to raise a discipline
%   issue, how to request supplies, marking policy)
% - A "Who to Call" quick reference card
%
% Format everything so it can be printed on 2 A4 pages.
% \end{lstlisting}
%
% \end{frame}
%
%
% %%%%%%%%%%%%%%%%%%%%%%%%%%%%%%%%%%%%%%%%%%%%%%%%%%%%%%%%%%%
% \begin{frame}[fragile]\frametitle{Analysing Survey or Feedback Data}
%
%     \begin{itemize}
%         \item \textbf{Situation:} You collected feedback forms, from students, parents, or clients, and need to make sense of 50 free-text responses.
%         \item \textbf{What would take hours} of reading, categorising, and writing is done in seconds. You validate and add context, the AI does the heavy lifting.
%     \end{itemize}
%
% \begin{lstlisting}[language=HTML]
% Prompt:
% Below are 50 responses from our annual parent
% satisfaction survey. Read all of them carefully.
%
% """
% [Paste all responses here]
% """
%
% Then provide:
% 1. Top 5 themes that appear most frequently
%    (with example quotes for each)
% 2. Top 3 specific complaints (most actionable)
% 3. Top 3 specific compliments (worth celebrating)
% 4. One surprising or unexpected insight
% 5. A short paragraph I can include in my report
%    summarising overall sentiment (positive/neutral/
%    negative split as a rough percentage)
% \end{lstlisting}
%
% \end{frame}

%%%%%%%%%%%%%%%%%%%%%%%%%%%%%%%%%%%%%%%%%%%%%%%%%%%%%%%%%%%
% [upgrade-deck: commented out 2026-10-24 | short sibling | left out of the 90-minute flipped seminar (full version keeps it)]
% \begin{frame}[fragile]\frametitle{Writing Personalised Student Feedback}
%
%     \begin{itemize}
%         \item \textbf{Situation:} You have 30 assignments to mark and need to write meaningful, individual feedback for each student, not a copy-paste.
%         \item \textbf{Key technique:} Give the AI the student's actual work and the rubric. It drafts feedback; you personalise and sign off.
% 		\item \textbf{Time saved:} 2--3 minutes per student instead of 10--15. Always review.% before sharing.
%
%     \end{itemize}
%
% \begin{lstlisting}[language=HTML]
% Prompt:
% You are a supportive teacher giving written feedback.
%
% Student name: Priya
% Assignment: Persuasive essay on school uniforms
% Grade level: 8th standard
%
% Rubric criteria: Argument clarity (5), Evidence (5), 
% Grammar (5), Structure (5)
%
% Student's essay: """[paste essay here]"""
%
% Write 3-4 sentences of encouraging, constructive 
% feedback. Mention one specific strength and one 
% clear area for improvement. Avoid generic praise.
% \end{lstlisting}
%
%
%
% \end{frame}


%%%%%%%%%%%%%%%%%%%%%%%%%%%%%%%%%%%%%%%%%%%%%%%%%%%%%%%%%%%
% [upgrade-deck: commented out 2026-10-24 | short sibling | left out of the 90-minute flipped seminar (full version keeps it)]
% \begin{frame}[fragile]\frametitle{Building a Personal ``Prompt Library''}
%
% \textbf{The highest-leverage habit you can build:}
%
% \medskip
% Once you find a prompt that works well for a recurring task, \textbf{save it as a reusable template}.
%
% \medskip
% Replace the \texttt{[PLACEHOLDERS]} each time. \textbf{Your prompt library becomes your personal AI productivity toolkit.}
%
% \begin{lstlisting}[language=HTML]
% Example saved prompt templates:
%
% [EMAIL_DELAY]
% Write a professional email informing [CLIENT] that 
% [PROJECT] is delayed by [DURATION] due to [REASON]. 
% Tone: apologetic but confident. Under 150 words.
%
% [LESSON_PLAN]  
% Create a 45-minute lesson plan for [TOPIC], Grade [X], 
% with prior knowledge of [Y]. Include warm-up, activity, 
% 5 practice problems, and a closing check question.
%
% [MEETING_SUMMARY]
% From these rough notes [NOTES], write: a 3-para summary, 
% an action table (Action | Owner | Deadline), and a 
% 4-line email to attendees.
% \end{lstlisting}
%
%
%
% \end{frame}


%%%%%%%%%%%%%%%%%%%%%%%%%%%%%%%%%%%%%%%%%%%%%%%%%%%%%%%%%%%
\begin{frame}[fragile]\frametitle{A Framework for Every Prompt: CRAFT}

\begin{center}
\Large \textbf{C R A F T}
\end{center}

\medskip
\begin{itemize}
    \item \textbf{C}ontext, who you are, what the situation is
    \item \textbf{R}ole, who the AI should act as
    \item \textbf{A}ction, exactly what you want it to do
    \item \textbf{F}ewshots, one or two examples of what ``good'' looks like
    \item \textbf{T}weaks, tone, length, format, and what to avoid
\end{itemize}

\medskip
\begin{center}
\textbf{Run through the five letters in your head before you type. That is the whole framework.}
\end{center}

\end{frame}

%%%%%%%%%%%%%%%%%%%%%%%%%%%%%%%%%%%%%%%%%%%%%%%%%%%%%%%%%%%
\begin{frame}[fragile]\frametitle{CRAFT: Context and Role}

\begin{itemize}
    \item \textbf{Context}, without it, the AI is guessing: what subject, what grade, what has already happened.
    \item \textbf{Role}, telling the AI who to act as changes its whole register, a ``supportive teacher'' writes differently than a ``strict examiner.''
\end{itemize}

\begin{lstlisting}[language=HTML]
Weak:  Explain photosynthesis.

CRAFT: You are a Grade 7 science teacher (Role).
       Your class has just finished a unit on plant
       cells and knows the word "chlorophyll" (Context).
\end{lstlisting}

\end{frame}

%%%%%%%%%%%%%%%%%%%%%%%%%%%%%%%%%%%%%%%%%%%%%%%%%%%%%%%%%%%
\begin{frame}[fragile]\frametitle{CRAFT: Action and Fewshots}

\begin{itemize}
    \item \textbf{Action}, one clear instruction. ``Help me with X'' is not an action, ``Write 5 quiz questions on X'' is.
    \item \textbf{Fewshots}, show, do not just tell. One example of the style or format you want is worth three sentences of description.
\end{itemize}

\begin{lstlisting}[language=HTML]
Action:   Write 5 short-answer questions on the
          water cycle.

Fewshot:  Match this style:
          "Q: What causes evaporation?
           A: Heat from the sun."
\end{lstlisting}

\end{frame}

%%%%%%%%%%%%%%%%%%%%%%%%%%%%%%%%%%%%%%%%%%%%%%%%%%%%%%%%%%%
\begin{frame}[fragile]\frametitle{CRAFT: Tweaks}

\begin{itemize}
    \item \textbf{Tweaks} are the small dials that turn a generic answer into a usable one:
    \begin{itemize}
        \item \textbf{Tone}, formal / friendly / encouraging
        \item \textbf{Length}, word count or ``under 5 sentences''
        \item \textbf{Format}, bullet list / table / email / paragraph
        \item \textbf{Constraints}, what to explicitly avoid
    \end{itemize}
    \item Add tweaks last, once Context, Role, Action, and Fewshots already give the AI enough to work with.
\end{itemize}

\begin{lstlisting}[language=HTML]
Tweaks: Keep it under 100 words, friendly tone,
        format as a bullet list, avoid technical jargon.
\end{lstlisting}

\end{frame}

%%%%%%%%%%%%%%%%%%%%%%%%%%%%%%%%%%%%%%%%%%%%%%%%%%%%%%%%%%%
\begin{frame}[fragile]\frametitle{CRAFT, All Together}

\begin{lstlisting}[language=HTML]
You are a supportive Grade 8 teacher (Role).      -- R
My class just finished a unit on fractions (C).   -- C

Write 5 short-answer practice questions on adding
fractions with different denominators (Action).   -- A

Match this style:
"Q: What is 1/2 + 1/4?  A: 3/4"                    -- F

Keep language simple, under 100 words total,
format as a numbered list.                         -- T
\end{lstlisting}

\end{frame}

%%%%%%%%%%%%%%%%%%%%%%%%%%%%%%%%%%%%%%%%%%%%%%%%%%%%%%%%%%%
\begin{frame}[fragile]\frametitle{CRAFT: The Takeaways}

\begin{itemize}
    \item Not every prompt needs all five letters spelled out explicitly, but running through them in your head catches what you would otherwise forget.
    \item \textbf{Always remember:} treat the output as a first draft, review it, and never paste real student names or confidential data into a public AI tool.
\end{itemize}

\end{frame}


%%%%%%%%%%%%%%%%%%%%%%%%%%%%%%%%%%%%%%%%%%%%%%%%%%%
\begin{frame}[fragile]\frametitle{Exercise: Weak Prompt to CRAFT Prompt}
\begin{itemize}
\item \textbf{2 minutes.} Run the weak prompt below with your own topic. Read the answer.
\item Rewrite it with all five CRAFT letters. Run it again.
\item Compare the two answers. Which one could you use tomorrow?
\end{itemize}
\begin{lstlisting}[language={}]
Weak prompt:
Make a worksheet on [TOPIC].
\end{lstlisting}
\end{frame}

\section[Flipped]{The Flipped Classroom, Briefly}

%%%%%%%%%%%%%%%%%%%%%%%%%%%%%%%%%%%%%%%%%%%%%%%%%%%
\begin{frame}[fragile]\frametitle{What Is a Flipped Classroom?}
\begin{itemize}
\item \textbf{Traditional:} the teacher explains in class. Students practise alone at home.
\item \textbf{Flipped:} students meet the content before class. Class time goes to discussion, problems and depth.
\item The teacher stays in the room for the hard part: questions, misconceptions, application.
\item The idea is old. AI makes the preparation cheap.
\end{itemize}
\end{frame}

%%%%%%%%%%%%%%%%%%%%%%%%%%%%%%%%%%%%%%%%%%%%%%%%%%%
\begin{frame}[fragile]\frametitle{Where AI Helps in a Flipped Classroom}
\begin{itemize}
\item \textbf{Before class, content:} reading, guiding questions, a diagram, a short quiz.
\item \textbf{In class, discussion:} questions, case studies, misconception probes, group tasks.
\item \textbf{After class, evaluation:} quiz with key, rubric, feedback drafts, exit ticket.
\item \textbf{Throughout, planning:} a minute-by-minute class plan.
\item You review, correct and decide. The five exercises follow this order.
\end{itemize}
\end{frame}

%%%%%%%%%%%%%%%%%%%%%%%%%%%%%%%%%%%%%%%%%%%%%%%%%%%
\begin{frame}[fragile]\frametitle{Choosing How Students Get the Content}
\begin{center}
\resizebox{\linewidth}{!}{%
\begin{tabular}{lcccc}
\toprule
 & Textbook & Own notes & Video playlist & LMS or MOOC \\
\midrule
Teacher effort & Low & Medium & Low & High \\
Matches your syllabus & Partly & Fully & Partly & Fully \\
Interactive & No & No & Little & Yes \\
Tracks students & No & No & No & Yes \\
\bottomrule
\end{tabular}}
\end{center}
\begin{itemize}
\item Routes combine well: one LMS unit can hold your notes, a video and a quiz.
\item Moodle and Open edX are two free LMS platforms. Hosting takes effort, so today uses plain notes.
\end{itemize}
\end{frame}

\section[Exercises]{Hands-On Exercises}

%%%%%%%%%%%%%%%%%%%%%%%%%%%%%%%%%%%%%%%%%%%%%%%%%%%
\begin{frame}[fragile]\frametitle{Get Ready for the Exercises}
\begin{itemize}
\item Open your AI assistant and sign in. Start a new chat.
\item Open one blank document. Paste every result there.
\item Choose \textbf{one topic} you teach: its subject, grade and a sample paragraph you like.
\item Use that same topic in every exercise.
\item Never type real student names or data into the tool.
\end{itemize}
\end{frame}

%%%%%%%%%%%%%%%%%%%%%%%%%%%%%%%%%%%%%%%%%%%%%%%%%%%
\begin{frame}[fragile]\frametitle{How the Exercises Run}
\begin{itemize}
\item Each exercise takes \textbf{9 minutes}: read 1, run 3, edit 3, share 2.
\item Replace every \texttt{[PLACEHOLDER]} with your own subject, grade and topic.
\item Work in pairs. Compare the outputs, not only the prompts.
\item Paste the best output into your document. It becomes your first lesson kit.
\item If an answer is poor, change the prompt and run it again. Do not accept the first draft.
\end{itemize}
\end{frame}

%%%%%%%%%%%%%%%%%%%%%%%%%%%%%%%%%%%%%%%%%%%%%%%%%%%
\begin{frame}[fragile]\frametitle{Exercise: Pre-Class Content}
\begin{itemize}
\item \textbf{Goal:} a pre-class reading with guiding questions for one topic.
\item \textbf{Step 1:} write the CRAFT prompt on the next slide for your topic.
\item \textbf{Step 2:} run it, then ask for an easier version.
\item \textbf{Step 3:} add one diagram, as shown two slides ahead.
\item \textbf{Result:} one page a student can read in 10 to 15 minutes.
\end{itemize}
\end{frame}

%%%%%%%%%%%%%%%%%%%%%%%%%%%%%%%%%%%%%%%%%%%%%%%%%%%
\begin{frame}[fragile]\frametitle{Exercise: Pre-Class Content, Starter Prompt}
\begin{itemize}
\item Change the bracketed parts. Keep the structure.
\end{itemize}
\begin{lstlisting}[language={}]
You are an experienced [GRADE] [SUBJECT] teacher (R).
My students know [PRIOR KNOWLEDGE] (C).
Write a pre-class reading on [TOPIC]: 400 words,
3 short headings, 1 worked example, 2 common
mistakes (A). Match the tone of this sample:
[PASTE A PARAGRAPH YOU LIKE] (F).
Simple words, short sentences. End with 3 guiding
questions to answer while reading (T).
\end{lstlisting}
\end{frame}

%%%%%%%%%%%%%%%%%%%%%%%%%%%%%%%%%%%%%%%%%%%%%%%%%%%
\begin{frame}[fragile]\frametitle{Exercise: Pre-Class Content, Diagram}
\begin{itemize}
\item SVG is text, so the AI can write it. Save the reply as \texttt{diagram.svg} and open it in a browser.
\item Ask for fixes in plain words. Other options: Mermaid text for flowcharts, or draw by hand from the AI's label list.
\end{itemize}
\begin{lstlisting}[language={}]
Draw an SVG, 400x200, of [PROCESS OR CONCEPT]:
list the parts and their positions here [PARTS],
arrows between them, every arrow and part labelled.
Output only the SVG code.
\end{lstlisting}
\end{frame}

%%%%%%%%%%%%%%%%%%%%%%%%%%%%%%%%%%%%%%%%%%%%%%%%%%%
\begin{frame}[fragile]\frametitle{Exercise: Pre-Class Content, Check}
\begin{itemize}
\item \textbf{Facts:} verify each number, date, name and formula against your textbook.
\item \textbf{Level:} read it as a student. Shorten any sentence you stumble over.
\item \textbf{Diagram:} check every label and arrow. Open it on a phone.
\item \textbf{Fit:} replace one example with one from your own town or school.
\end{itemize}
\end{frame}

%%%%%%%%%%%%%%%%%%%%%%%%%%%%%%%%%%%%%%%%%%%%%%%%%%%
\begin{frame}[fragile]\frametitle{Exercise: In-Class Discussion}
\begin{itemize}
\item \textbf{Goal:} a 25-minute in-class activity that builds on the reading.
\item \textbf{Step 1:} run the prompt on the next slide with the same topic.
\item \textbf{Step 2:} pick the two best questions and one group task.
\item \textbf{Step 3:} try one more format from the slide after.
\item \textbf{Result:} questions and a task you can use in the next class.
\end{itemize}
\end{frame}

%%%%%%%%%%%%%%%%%%%%%%%%%%%%%%%%%%%%%%%%%%%%%%%%%%%
\begin{frame}[fragile]\frametitle{Exercise: In-Class Discussion, Starter Prompt}
\begin{itemize}
\item Paste your pre-class reading below the prompt, so the AI knows what students have read.
\end{itemize}
\begin{lstlisting}[language={}]
You are a [GRADE] [SUBJECT] teacher running a flipped
class (R). Students read the notes below at home (C).
Write: 5 discussion questions, from factual to
open-ended; 3 common misconceptions, each with one
probing question; 1 group task for 4 students that
takes 15 minutes (A).
No question may be answered by copying from the
notes (T).
Notes: [PASTE YOUR READING]
\end{lstlisting}
\end{frame}

%%%%%%%%%%%%%%%%%%%%%%%%%%%%%%%%%%%%%%%%%%%%%%%%%%%
\begin{frame}[fragile]\frametitle{Exercise: In-Class Discussion, Other Formats}
\begin{itemize}
\item \textbf{Case study:} a short real-life scenario with 3 questions.
\item \textbf{Debate:} two positions on the topic, with a role card for each group.
\item \textbf{Think-pair-share:} one question, 1 minute alone, 2 in pairs, then a class answer.
\item \textbf{Devil's advocate:} ask the AI to argue a wrong view. Students find the flaw.
\item Ask for the format by name, plus the group size and the time.
\end{itemize}
\end{frame}

%%%%%%%%%%%%%%%%%%%%%%%%%%%%%%%%%%%%%%%%%%%%%%%%%%%
\begin{frame}[fragile]\frametitle{Exercise: In-Class Discussion, Check}
\begin{itemize}
\item \textbf{Open ended:} would two students give different good answers?
\item \textbf{Misconceptions:} are they real mistakes your students make?
\item \textbf{Time:} can the group task finish in 15 minutes with your classroom materials?
\item \textbf{Safety:} is every example and role respectful of your students' backgrounds?
\end{itemize}
\end{frame}

%%%%%%%%%%%%%%%%%%%%%%%%%%%%%%%%%%%%%%%%%%%%%%%%%%%
\begin{frame}[fragile]\frametitle{Exercise: Evaluation}
\begin{itemize}
\item \textbf{Goal:} a short quiz with an answer key, and a rubric for one task.
\item \textbf{Step 1:} generate five questions with answers and reasons.
\item \textbf{Step 2:} generate a rubric for one assignment on the same topic.
\item \textbf{Step 3:} draft feedback on a \emph{made-up} student answer.
\item \textbf{Result:} a check quiz and a rubric ready to print.
\end{itemize}
\end{frame}

%%%%%%%%%%%%%%%%%%%%%%%%%%%%%%%%%%%%%%%%%%%%%%%%%%%
\begin{frame}[fragile]\frametitle{Exercise: Evaluation, Quiz Prompt}
\begin{itemize}
\item Ask for wrong options built from real misconceptions. Keep the key apart from the student copy.
\end{itemize}
\begin{lstlisting}[language={}]
You are a [GRADE] [SUBJECT] teacher (R).
Write 5 multiple-choice questions on [TOPIC],
4 options each, mixing recall and application (A).
Make each wrong option a common misconception (T).
Return a table: Question | Correct | Why it is right |
Why each wrong option is wrong (T).
\end{lstlisting}
\end{frame}

%%%%%%%%%%%%%%%%%%%%%%%%%%%%%%%%%%%%%%%%%%%%%%%%%%%
\begin{frame}[fragile]\frametitle{Exercise: Evaluation, Rubric and Feedback}
\begin{itemize}
\item Run the rubric prompt first. Then paste a made-up answer into the second prompt.
\item The AI drafts the feedback. You decide the mark.
\end{itemize}
\begin{lstlisting}[language={}]
Prompt 1:
Create a rubric for [ASSIGNMENT] for [GRADE] with
4 criteria and 3 levels. One short descriptor per
cell, as a table.

Prompt 2:
Using this rubric, write 3 sentences of feedback:
one strength, one improvement, one next step.
Answer: [PASTE A MADE-UP STUDENT ANSWER]
\end{lstlisting}
\end{frame}

%%%%%%%%%%%%%%%%%%%%%%%%%%%%%%%%%%%%%%%%%%%%%%%%%%%
\begin{frame}[fragile]\frametitle{Exercise: Evaluation, Check}
\begin{itemize}
\item \textbf{Answer key:} answer every question yourself. Does the key match?
\item \textbf{Wrong options:} would a real student pick them?
\item \textbf{Rubric:} are the levels clear enough that two teachers would mark alike?
\item \textbf{Feedback:} does it name something specific in the answer, not generic praise?
\item \textbf{Privacy:} no real names, marks or health details in any prompt.
\end{itemize}
\end{frame}

%%%%%%%%%%%%%%%%%%%%%%%%%%%%%%%%%%%%%%%%%%%%%%%%%%%
\begin{frame}[fragile]\frametitle{Exercise: Plan One Flipped Lesson}
\begin{itemize}
\item \textbf{Goal:} a one-page, minute-by-minute plan for one flipped class.
\item \textbf{Step 1:} paste your pre-class task and your best discussion activity into the prompt.
\item \textbf{Step 2:} ask what could go wrong, and what to do for unprepared students.
\item \textbf{Result:} a plan you could teach next week.
\end{itemize}
\end{frame}

%%%%%%%%%%%%%%%%%%%%%%%%%%%%%%%%%%%%%%%%%%%%%%%%%%%
\begin{frame}[fragile]\frametitle{Exercise: Plan One Flipped Lesson, Starter Prompt}
\begin{itemize}
\item Follow with: ``What could go wrong, and what is my backup?''
\end{itemize}
\begin{lstlisting}[language={}]
You are a [GRADE] [SUBJECT] teacher planning a
flipped lesson (R). Students have done this
pre-class task: [PASTE TASK] (C).
Plan a 45-minute class: 5-minute check quiz,
10 minutes on what students found confusing,
25-minute activity [PASTE ACTIVITY], 5-minute
close with one exit question (A).
Return a table: Minutes | Teacher does |
Students do. Use only a board and paper (T).
\end{lstlisting}
\end{frame}

%%%%%%%%%%%%%%%%%%%%%%%%%%%%%%%%%%%%%%%%%%%%%%%%%%%
\begin{frame}[fragile]\frametitle{Exercise: Plan One Flipped Lesson, Check}
\begin{itemize}
\item \textbf{Timing:} do the minutes add up to 45? Is the transition time realistic?
\item \textbf{Materials:} can you run it with what your classroom has?
\item \textbf{Backup:} what happens when half the class skipped the reading?
\item \textbf{Fit:} does the plan match your students' real level and class size?
\end{itemize}
\end{frame}

%%%%%%%%%%%%%%%%%%%%%%%%%%%%%%%%%%%%%%%%%%%%%%%%%%%
\begin{frame}[fragile]\frametitle{Exercise: Fellowship Project Ideas}
\begin{itemize}
\item \textbf{Goal:} three topic options and one draft problem statement for your fellowship project.
\item \textbf{Step 1:} describe your real setting to the AI and ask for topic options.
\item \textbf{Step 2:} answer its questions, and choose one topic.
\item \textbf{Step 3:} write the problem statement yourself, then ask the AI to critique it.
\item \textbf{Result:} a topic you can take into the proposal-writing session.
\end{itemize}
\end{frame}

%%%%%%%%%%%%%%%%%%%%%%%%%%%%%%%%%%%%%%%%%%%%%%%%%%%
\begin{frame}[fragile]\frametitle{Exercise: Fellowship Project Ideas, Topic Prompt}
\begin{itemize}
\item Be honest about your constraints: devices, internet, class size, time.
\end{itemize}
\begin{lstlisting}[language={}]
You are an education researcher mentoring a school
teacher (R). I teach [SUBJECT] to [GRADE] in a
[TYPE OF SCHOOL]. My constraints: [DEVICES, INTERNET,
CLASS SIZE, TIME] (C).
Suggest 5 project topics on using a flipped
classroom with AI in my setting. For each: the
problem it solves and one measurable outcome (A).
Then ask me 5 questions to choose the best one (T).
\end{lstlisting}
\end{frame}

%%%%%%%%%%%%%%%%%%%%%%%%%%%%%%%%%%%%%%%%%%%%%%%%%%%
\begin{frame}[fragile]\frametitle{Exercise: Fellowship Project Ideas, Critique Prompt}
\begin{itemize}
\item Write 3 to 4 sentences yourself first. The AI critiques; it does not write your proposal.
\end{itemize}
\begin{lstlisting}[language={}]
Here is my draft problem statement:
[PASTE YOUR 3 TO 4 SENTENCES]
List its weaknesses: vague words, no measure,
too big for one year, missing baseline.
Then list 5 questions a reviewer would ask me.
Do not rewrite it.
\end{lstlisting}
\end{frame}

%%%%%%%%%%%%%%%%%%%%%%%%%%%%%%%%%%%%%%%%%%%%%%%%%%%
\begin{frame}[fragile]\frametitle{Exercise: Fellowship Project Ideas, Check}
\begin{itemize}
\item \textbf{Your class:} does the topic fit your real students, school and year?
\item \textbf{Measurable:} could you collect the outcome with tools you have?
\item \textbf{References:} the AI may invent papers and links. Search for each one yourself.
\item \textbf{Ownership:} the proposal text is yours. Use the AI as a mentor.
\end{itemize}
\end{frame}

\section[Checks]{Use It Responsibly}

%%%%%%%%%%%%%%%%%%%%%%%%%%%%%%%%%%%%%%%%%%%%%%%%%%%
\begin{frame}[fragile]\frametitle{Check Before You Use It}
\begin{itemize}
\item \textbf{Facts:} verify every number, date, name and formula.
\item \textbf{Level:} read it as a student would.
\item \textbf{Bias:} look at examples, names and tone, especially in feedback.
\item \textbf{Ownership:} edit until it sounds like you.
\item \textbf{Privacy:} no real student data in a public tool.
\end{itemize}
\end{frame}

%%%%%%%%%%%%%%%%%%%%%%%%%%%%%%%%%%%%%%%%%%%%%%%%%%%
\begin{frame}[fragile]\frametitle{Student Data: The One Non-Negotiable}
\begin{itemize}
\item \textbf{Never paste} student names, marks, health details or other personal data into a public AI tool.
\item Use an approved school tool for real student records, if your school has one.
\item Anonymise first: ``a Grade 7 student who struggles with fractions'' works as well as a name.
\end{itemize}
\begin{alertblock}{Rule of Thumb}
If you would not write it on a sticky note left on a public desk, do not paste it into a public AI chat.
\end{alertblock}
\end{frame}

%%%%%%%%%%%%%%%%%%%%%%%%%%%%%%%%%%%%%%%%%%%%%%%%%%%
\begin{frame}[fragile]\frametitle{Quick Check: Spot the Problem}
\begin{itemize}
\item A colleague pastes this into a public AI chatbot to draft feedback faster.
\item What is wrong here, and what is the one-line fix?
\end{itemize}
\begin{lstlisting}[language={}]
Write encouraging feedback for Priya Sharma, roll no. 23,
who has been diagnosed with ADHD and is struggling with
her Grade 8 persuasive essay. Her essay: [pastes full text]
\end{lstlisting}
\end{frame}

%%%%%%%%%%%%%%%%%%%%%%%%%%%%%%%%%%%%%%%%%%%%%%%%%%%
\begin{frame}[fragile]\frametitle{Quick Check: Answer}
\begin{itemize}
\item \textbf{The problem:} a full name, a roll number and a health diagnosis, all sent to a public tool.
\item \textbf{The fix:} drop the name, roll number and diagnosis. Keep only what the task needs: ``a Grade 8 student who needs extra encouragement and structure'', plus the essay text.
\item Same quality of feedback, no exposure.
\end{itemize}
\end{frame}

%%%%%%%%%%%%%%%%%%%%%%%%%%%%%%%%%%%%%%%%%%%%%%%%%%%
\begin{frame}[fragile]\frametitle{Four Things That Go Wrong}
\begin{itemize}
\item \textbf{Hallucination:} confident wrong answers. Check everything that reaches students.
\item \textbf{Academic integrity:} students use AI to write essays. Redesign tasks: in-class writing, oral defence, drafts.
\item \textbf{Over-reliance:} the first draft is a starting point, for you and for students.
\item \textbf{Bias:} models inherit the biases of their training text. Review feedback and grading help with that in mind.
\end{itemize}
\end{frame}

%%%%%%%%%%%%%%%%%%%%%%%%%%%%%%%%%%%%%%%%%%%%%%%%%%%
\begin{frame}[fragile]\frametitle{Your Next Seven Days}
\begin{itemize}
\item Pick \textbf{one} output from today: the reading, the discussion task or the quiz.
\item Use it with a real class this week. Note what worked and what confused students.
\item Save the prompt that worked, with placeholders, in your own prompt library.
\item Bring one result to your fellowship project plan.
\end{itemize}
\end{frame}

%%%%%%%%%%%%%%%%%%%%%%%%%%%%%%%%%%%%%%%%%%%%%%%%%%%
\begin{frame}[fragile]\frametitle{Takeaways}
\begin{itemize}
\item Flipping moves content before class and gives class time to discussion and application.
\item AI drafts content, discussion material, quizzes and plans. You check, correct and own them.
\item A CRAFT prompt turns a generic answer into a usable one.
\item Protect student data. Start small, and share what you learn.
\end{itemize}
\end{frame}
